\begin{abstract}
Federated link prediction is difficult to deploy privately: interaction edges
remain distributed across clients, yet a useful predictor must exploit graph
structure and answer candidate-pair queries without rereading the private
graph. A second difficulty is less visible but equally important: under edge
differential privacy (DP), a structural channel can improve one target graph
and harm another, so unconditional deployment can be worse than a predictor
using only public node descriptors. We introduce \method{}, an
architecture-agnostic policy that combines an inference-closed edge-DP
structural learner with a private, target-domain utility certificate.
\method{} activates the structural branch only when a one-sided lower
certificate exceeds a registered material-gain threshold; otherwise it falls
back to the public branch. We prove that no training-only selector can provide
a distribution-free nontrivial no-harm guarantee, establish complete
training-and-certification transcript privacy, and derive a finite-population
no-harm theorem. Sufficient and necessary certification counts match in their
principal dependence on effect gap, privacy budget, and record dependence. In
a preregistered one-time evaluation on six social and blog networks with five
seeds, the primary visible-message mechanism has composed privacy
$(\eps=5.664,\delta=2\!\times\!10^{-6})$. It activates 15 of 30
dataset-seed cells across three datasets, observes no material false
activation, and obtains mean disjoint-holdout pairwise gain $0.0854$ over the
public branch. An always-structural policy is harmful in 10 of 30 cells and
achieves lower mean gain. The result supports certified fallback, rather than
universal structural superiority, as a defensible deployment principle for
federated edge-private link prediction.
\end{abstract}

\begin{IEEEkeywords}
Differential privacy, federated graph learning, link prediction, private
certification, privacy-utility trade-off.
\end{IEEEkeywords}
